\documentclass{article}

\usepackage{amsmath,amssymb}
\usepackage{xurl}
\usepackage[colorlinks=true,linkcolor=blue,citecolor=blue,urlcolor=blue]{hyperref}

% Shared commands for notes in docs/paper-gaps/.

\DeclareUnicodeCharacter{2080}{\ifmmode{}_{0}\else\textsubscript{0}\fi}
\DeclareUnicodeCharacter{2081}{\ifmmode{}_{1}\else\textsubscript{1}\fi}
\DeclareUnicodeCharacter{2082}{\ifmmode{}_{2}\else\textsubscript{2}\fi}
\DeclareUnicodeCharacter{2099}{\ifmmode{}_{n}\else\textsubscript{n}\fi}

\newcommand{\Lean}{\textsc{Lean}}
\ExplSyntaxOn
\NewDocumentCommand{\leanid}{m}{
  \ifmmode
    \textsf{\detokenize{#1}}
  \else
    \group_begin:
    \urlstyle{sf}
    \tl_set:Nn \l_tmpa_tl {#1}
    \tl_replace_all:Nnn \l_tmpa_tl {\_} {_}
    \exp_args:NV \path \l_tmpa_tl
    \group_end:
  \fi
}
\ExplSyntaxOff
\providecommand{\tr}{\operatorname{tr}}
\newcommand{\tnleanIssueBase}{https://github.com/LionSR/TNLean/issues/}
\newcommand{\tnleanPrBase}{https://github.com/LionSR/TNLean/pull/}
\newcommand{\tnleanDocBase}{https://sirui-lu.com/TNLean/docs/}
\newcommand{\ghissue}[1]{\href{\tnleanIssueBase#1}{GitHub issue \##1}}
\newcommand{\ghpr}[1]{\href{\tnleanPrBase#1}{PR \##1}}
\newcommand{\doclink}[1]{\href{\tnleanDocBase#1.html}{\path{#1}}}

% Dirac notation and the identity operator, as in blueprint/src/macros/common.tex.
% \providecommand keeps note-local definitions authoritative.
\usepackage{dsfont}
\providecommand{\ket}[1]{|#1\rangle}
\providecommand{\bra}[1]{\langle #1|}
\providecommand{\braket}[2]{\langle #1 | #2 \rangle}
\providecommand{\ketbra}[2]{|#1\rangle\!\langle #2|}
\providecommand{\Id}{\mathds{1}}
\providecommand{\one}{\mathds{1}}
\providecommand{\C}{\mathbb{C}}
\providecommand{\MN}[1]{M_{#1}(\C)}
\providecommand{\MD}{M_D(\C)}

% External referencing for paper sources.  The build helper writes sanitized
% auxiliary files under build/paper-aux/ and excludes duplicated source labels.
\usepackage{xr}

% Import a generated paper auxiliary file with a caller-supplied label prefix.
% The candidate paths support builds from docs/paper-gaps/, the repository root,
% and build/paper-aux/ itself.
\newcommand{\importpaperaux}[2]{%
  \IfFileExists{../../build/paper-aux/#2.aux}{%
    \externaldocument[#1]{../../build/paper-aux/#2}%
  }{%
    \IfFileExists{build/paper-aux/#2.aux}{%
      \externaldocument[#1]{build/paper-aux/#2}%
    }{%
      \IfFileExists{#2.aux}{%
        \externaldocument[#1]{#2}%
      }{}%
    }%
  }%
}

% General external reference macros
\newcommand{\paperref}[2]{\ref{#1-#2}}
\newcommand{\papereqref}[2]{(\ref{#1-#2})}

% CPSV16 source (1606.00608 / MPDO-22-12-17-2.tex) external document and shortcuts
\importpaperaux{cpsv16-}{1606.00608/MPDO-22-12-17-2-tnlean}
\newcommand{\cpsvref}[1]{\paperref{cpsv16}{#1}}
\newcommand{\cpsveqref}[1]{\papereqref{cpsv16}{#1}}

% Verdict marker: \gapnote{<kind>}{<status>}. Kind is the policy
% classification (clarification, local-correction, scope-restriction,
% unfaithful, false-source, open-gap); status is open, wip (elimination
% underway), resolved, or historical. Machine-read by the site index;
% prints a verdict line.
\newcommand{\gapnote}[2]{\par\noindent\textbf{Verdict:} #1 (#2).\par}


\title{Orientation of the Three-Swap Ancilla Circuit}
\author{}
\date{2026-08-31}

\begin{document}
\maketitle
\gapnote{false-source}{resolved}

\section{The printed circuit}

Proposition \texttt{prop:U1-U2-equiv-ancillatrick} of
Ref.~\cite{Cirac2017MPU} claims that, after adjoining one identity ancilla per
site, conjugating $\widetilde U_1$ by the three swap layers in
Figure~\texttt{fig:TR-ancilla} gives $\widetilde U_2$. The layers above
$\widetilde U_1$, read from the endpoint outwards, are labelled
$S^{1,a}_{n,n+1}$, $S^{1,2}_{n,n}$, and $S^{2,a}_{n,n}$.

Let $R$ denote cyclic rotation of configurations, with
$(R\sigma)(n)=\sigma(n+1)$. On triples $(x,y,a)$ of the first species, second
species, and ancilla configurations, the three literal swap labels act as
\begin{align}
    A(x,y,a)
     & =(Ra,y,R^{-1}x),
    \label{eq:mpu_ancilla_orientation_crossed_swap}  \\
    B(x,y,a)
     & =(y,x,a),
    \label{eq:mpu_ancilla_orientation_physical_swap} \\
    D(x,y,a)
     & =(x,a,y).
    \label{eq:mpu_ancilla_orientation_ancilla_swap}
\end{align}
Here $A$ is the global product of the swaps between species $1$ at site $n$
and the ancilla at site $n+1$. The first output therefore receives the
next-site ancilla, while the ancilla output receives the previous-site first
species.

\section{The orientation mismatch}

Let $C=D\circ B\circ A$ be the stack above $\widetilde U_1$. Direct
substitution into (\ref{eq:mpu_ancilla_orientation_crossed_swap})--
(\ref{eq:mpu_ancilla_orientation_ancilla_swap}) gives
\begin{align}
    C(x,y,a)
     & =(y,R^{-1}x,Ra),
    \label{eq:mpu_ancilla_orientation_stack} \\
    C^{-1}(x,y,a)
     & =(Ry,x,R^{-1}a).
    \label{eq:mpu_ancilla_orientation_stack_inverse}
\end{align}
Consequently, conjugating the physical swap gives
\begin{align}
    C^{-1}BC(x,y,a)
     & =(Ry,R^{-1}x,a).
    \label{eq:mpu_ancilla_orientation_literal_endpoint}
\end{align}

The paper defines $T$ by
$T\lvert n_1,\ldots,n_N\rangle
    =\lvert n_N,n_1,\ldots,n_{N-1}\rangle$. Thus $T$ sends a ket configuration
$	au$ to $R^{-1}\tau$, and $T^\dagger$ sends it to $R\tau$. It follows that
the two swap-transformed shifts act by
\begin{align}
    \widetilde U_2(x,y,a)
     & =(R^{-1}y,Rx,a),
    \label{eq:mpu_ancilla_orientation_tilde_u2} \\
    \widetilde U_3(x,y,a)
     & =(Ry,R^{-1}x,a).
    \label{eq:mpu_ancilla_orientation_tilde_u3}
\end{align}
Comparison of (\ref{eq:mpu_ancilla_orientation_literal_endpoint}) with
(\ref{eq:mpu_ancilla_orientation_tilde_u2})--
(\ref{eq:mpu_ancilla_orientation_tilde_u3}) shows that the literal circuit
labels always agree with $\widetilde U_3$. For $N\geq3$ and $d\geq2$, they do
not in general agree with $\widetilde U_2$: one may choose a configuration
whose values at the preceding and following sites differ, so
$R\sigma\neq R^{-1}\sigma$. By contrast, $R=R^{-1}$ when $N\leq2$, and the
configuration space has at most one element when $d\leq1$. The two endpoint
permutations therefore coincide in these degenerate cases.

\section{Local correction and formal record}

Reversing the crossed swap replaces $A$ by
\begin{align}
    A_{\mathrm{rev}}(x,y,a)
     & =(R^{-1}a,y,Rx).
    \label{eq:mpu_ancilla_orientation_reversed_crossed_swap}
\end{align}
This is the swap between species $1$ at site $n+1$ and the ancilla at site
$n$, rather than the literal $S^{1,a}_{n,n+1}$ label. Using
$A_{\mathrm{rev}}$ in the same three-layer stack changes
(\ref{eq:mpu_ancilla_orientation_literal_endpoint}) to
\begin{align}
    (x,y,a)
     & \longmapsto(R^{-1}y,Rx,a),
    \label{eq:mpu_ancilla_orientation_corrected_endpoint}
\end{align}
which is the $\widetilde U_2$ endpoint in
(\ref{eq:mpu_ancilla_orientation_tilde_u2}). Thus reversing the orientation
of the crossed labels is a local correction of the diagram; it does not alter
the proposition's intended endpoint.

The formal coordinate calculation keeps the literal printed labels and proves
(\ref{eq:mpu_ancilla_orientation_literal_endpoint}). It asserts only this
permutation identity and does not identify its permutation matrix with either
ancilla-enlarged MPO or use it as a dependency of the
$\widetilde U_1$--$\widetilde U_2$ path proposition.
\footnote{Formal declarations:
    \leanid{MPOTensor.shiftAncillaCounterShift} and
    \leanid{MPOTensor.permMatrix_shiftAncillaThreeSwap_endpoint} in
    \doclink{TNLean/MPS/MPU/Examples/ShiftSymmetryPaths}.}

\section{Verdict}

The three swap labels in Figure~\texttt{fig:TR-ancilla} agree with
$\widetilde U_3$ under the shift convention stated in the same paper. They
also agree with $\widetilde U_2$ in the degenerate cases $N\leq2$ or $d\leq1$,
but the figure's claimed $\widetilde U_2$ endpoint is false as printed on
nondegenerate periodic chains. Reversing the orientation of the crossed
species--ancilla swaps gives the intended endpoint, so the discrepancy has a
resolved local correction.

\bibliographystyle{alpha}
\bibliography{references}

\end{document}